\documentclass[10pt,a4paper]{amsart}
\usepackage[margin=19mm]{geometry}
\usepackage{amsmath,amssymb,mathtools}
\usepackage[scr=rsfs,cal=euler]{mathalfa}
\usepackage{xfrac}
\usepackage[unicode,hidelinks,bookmarks=false]{hyperref}
\newcommand{\ie}{i.e.\ }
\newcommand{\cf}{cf.\ }
\newcommand{\resp}{respectively\ }
\newcommand{\Subsection}{\subsection}
\providecommand{\nobreakdash}{\nobreak\hbox{-}}
\setlength{\emergencystretch}{2em}
\multlinegap0pt
\usepackage{bm}
\usepackage{bbm}
\usepackage{dsfont}
\usepackage{xspace}
\usepackage{cancel}
\usepackage{mathtools}
\usepackage{amsthm}
\makeatletter
\@ifundefined{Lemma}{\newtheorem{Lemma}{Lemma}}{}
\makeatother
\title[On a class of Cauchy problems with applications in nonlinear]{On a class of Cauchy problems with applications in nonlinear partial differential equations}
\author{Feida Jiang, Neil S. Trudinger, Qiao-Qiao Xu}
\date{Source: arXiv:2601.05550v1 (CC BY 4.0)}
\begin{document}
\maketitle

\noindent\emph{Source and licence.} On a class of Cauchy problems with applications in nonlinear partial differential equations. Version arXiv:2601.05550v1 (CC BY 4.0), distributed under the Creative Commons Attribution 4.0 International licence, as recorded in the arXiv licence metadata for this exact version and shown on the abstract page https://arxiv.org/abs/2601.05550v1. Full rights notice: licenses/arxiv-2601.05550-CC-BY-4.0.txt.\par\smallskip

\noindent\textit{Editorial scope.} Counted unit: the Lemma whose source label is \texttt{Lem 3.1} in the article (source0.tex lines 540--547, source label \texttt{Lem 3.1}), delivered together with the article's own complete original proof of it (source0.tex lines 548--610, one \texttt{proof} environment of 63 lines). No mathematics is written, repaired or re-derived: statement and proof are the article's own text line for line. Required context retained uncounted: 1.1 (lines 103--110); KO (lines 128--130); 3.1 (lines 533--536). Every definition, display and label of those lines is retained unchanged. Omitted from this package and not redistributed: 1--102, 111--127, 131--532, 537--539, 611--1685. Nothing inside a retained excerpt is omitted or altered: every retained body is delivered exactly as the article's lines are written, so no omission receipt of any kind accompanies this package. Citations: the retained excerpts carry no citation command at all, so no bibliography entry of the article is transcribed and no reference is redistributed. Reference adaptation: none was needed and none was made; every reference inside the delivered unit resolves inside it, so no unresolved reference or citation is produced. Printed slips kept literal: no printed word, symbol, hypothesis or number of the counted statement or of its proof was altered. Layout and class: the article's own class is \texttt{amsart} with packages including amsmath, url, graphicx, amssymb, color, pstricks-add, manfnt, amsthm, amsfonts, mathrsfs, multirow, cite, enumerate, hyphenat; the wrapper is the lane's pinned \texttt{amsart} editorial wrapper at 10pt on A4, which restates the theorem-like environments the retained text needs and adds the article's own macro definitions that the retained text needs (none), with extra wrapper packages (bm, bbm, dsfont, xspace, cancel, mathtools). The delivered numbering is the unit's own. Assets: no retained span contains a figure environment, a graphics-inclusion command, a PDF-page inclusion, a table or a TikZ picture; no image file is copied into this package, the package declares no asset dependency and no image file is redistributed with it. Licence: version arXiv:2601.05550v1 of this article is distributed under the Creative Commons Attribution 4.0 International licence, as recorded in the arXiv licence metadata for this exact version and shown on the abstract page https://arxiv.org/abs/2601.05550v1. A full rights notice is delivered at licenses/arxiv-2601.05550-CC-BY-4.0.txt. Characters: every retained excerpt is pure ASCII, so the delivered engine, XeLaTeX with its default Latin Modern text font, reports no missing character.
	\begin{equation}\label{1.1}
	\left\{
	\begin{aligned}
	&v'(r)= C  r^{-q}  \left(  \int_{0}^{r} s^{\tau-1} g (v(s))ds \right )^{\frac{1}{\theta}}, \quad r\in (0, +\infty), \\
	&v(0)=a,
	\end{aligned}
	\right.
	\end{equation}

	\begin{equation}\label{KO}
	\int^{+\infty} \left ( \int_0^t g(s)ds \right )^{-\frac{1}{\theta+1}} dt = + \infty,
	\end{equation}

\begin{equation}\label{3.1}
v^{\prime\prime}(r)\left( v^{\prime}(r)\right)^{\theta -1}+\frac{q}{r}\left( v^{\prime}(r)\right)^{\theta}
=\frac{C^\theta}{\theta} r^{\tau - q \theta -1} g(v(r)).
\end{equation}

\begin{Lemma}\label{Lem 3.1}
Assume that $g: \mathbb{R} \to (0,+\infty)$ is a non-decreasing continuous function, the constants $q\ge 0$, $\theta> 0$, $C>0$ and $\tau > \max \{\theta (q-1), 0\}$, and $v$ is a solution of the Cauchy problem \eqref{1.1} with $v'(0)=0$ for all $r\ge 0$ and for some initial value $a$.
If condition \eqref{KO} does not hold, then there exists a large positive constant $R_0$ such that
\begin{equation}\label{claim}
v^{\prime\prime}(r)\left( v^{\prime}(r)\right)^{\theta -1} > \frac{C^\theta}{2\theta} r^{\tau - q \theta -1} g(v(r)),
\end{equation}
for $r>R_0$.
\end{Lemma}

\begin{proof}[Proof of Lemma \ref{Lem 3.1}]
Letting
$\tilde g(t)=\left( \int_{0}^{t}g (s)ds\right)^{-\frac{1}{\theta+1}}$, since condition \eqref{KO} deos not hold, we have  $\int^{\infty}\tilde g(t) dt<\infty$. Hence, $\int_s^{\infty}\tilde g(t)dt\to 0$ as $s\to \infty$. In addition, by observing that $\tilde g(t)$ is decreasing in $(0, \infty)$, we have
\begin{equation*}
	0<t \tilde g(t)\leq 2\int_{\frac{t}{2}}^{t}\tilde g(s)ds
	<2\int_{\frac{t}{2}}^{\infty}\tilde g(s)ds \to 0, \quad {\rm as} \quad t\to \infty,
\end{equation*}
which shows that
\begin{equation*}
	t \tilde g(t)=t \left( \int_{0}^{t}g (s)ds\right)^{-\frac{1}{\theta+1}}\to 0,\quad {\rm as} \quad t\to \infty.
\end{equation*}
Combining this with the fact that $g$ is non-decreasing, we have
\begin{equation*}
	\left(\frac{ g^{\frac{1}{\theta}}(t)}{t} \right)^\theta =\frac{t g(t)}{t^{\theta+1}}
	\geq\frac{\int_{0}^{t}g(s)ds}{t^{\theta+1}}=(t\tilde g(t))^{-(\theta+1)}\to \infty, \quad {\rm as} \quad t\to \infty.
\end{equation*}
Therefore there exists a constant $t_1 > a$ such that 	
\begin{equation}\label{3.6'}
	g^{\frac{1}{\theta}}(t)>t-a,\quad \forall \,t>t_1.
\end{equation}
Since $v^{\prime}(r)>0$ for $r>0$, we have $v(r)>v(0)=a$ for $r>0$.
Since $g$ is non-decreasing and $v(0)=a$, we have from \eqref{1.1} that 
$$v^\prime(r) \ge Cr^{-q}g^{\frac{1}{\theta}}(a)\left (\int_0^r s^{\tau-1} ds \right)^{\frac{1}{\theta}} = C\left (\frac{1}{\tau}\right)^{\frac{1}{\theta}}r^{\frac{\tau}{\theta}-q} g^{\frac{1}{\theta}}(a).$$
Integrating from $0$ to $r$ and using $v(0)=a$ again, we have
$$v(r)\geq C\left (\frac{1}{\tau}\right)^{\frac{1}{\theta}} \frac{\theta}{\tau - q \theta + \theta}g^{\frac{1}{\theta}}(a) r^{\frac{\tau}{\theta}-q+1}  +a, \quad   r>0.$$
Since $\tau > \max \{\theta (q-1), 0\}$, from the above inequality, there exists $\tilde r>0$ such that
$$v(r)>t_1, \quad {\forall} r>\tilde r.$$
Therefore,  from \eqref{3.6'}, we have
\begin{equation}\label{3.10}
	g^{\frac{1}{\theta}}(v(r))>v(r)-a, \quad   r>\tilde r.
\end{equation}
Since $v'(r)>0$ for $r>0$, from \eqref{3.1}, we have
\begin{equation*}
v^{\prime\prime}(r)\left( v^{\prime}(r)\right)^{\theta - 1}
\le \frac{1}{\theta} C^\theta r^{\tau - q\theta -1} g(v(r)), \quad r>0.
\end{equation*}
Multiplying this by $(\theta+1)v'(r)$, we have
\begin{equation*}
(\theta+1)v^{\prime\prime}(r)\left( v^{\prime}(r)\right)^{\theta}
\le \frac{\theta + 1}{\theta} C^\theta r^{\tau - q\theta -1} g(v(r)) v^{\prime}(r), \quad r>0.
\end{equation*}
Integrating this from $0$ to $r(>\tilde r)$, we get
\begin{equation}\label{3.11}
\begin{aligned}
		\left(v^{\prime}(r)\right)^{\theta + 1} & \le \frac{\theta + 1}{\theta} C^\theta r^{\tau - q\theta -1} g(v(r))(v(r)-a) \\
		& < \frac{\theta + 1}{\theta} C^\theta r^{\tau - q\theta -1} g^{1+\frac{1}{\theta}}(v(r)), \quad r> \tilde r,
\end{aligned}
\end{equation}
where $v^\prime(0)=0$ and the monotonicity of $g$ are used in the first inequality, and \eqref{3.10} is used in the second inequality.
By direct calculation, we obtain from \eqref{3.11} that
\begin{equation}\label{3.12}
\frac{q}{r}\left( v^{\prime}(r)\right)^{\theta}< \left[ q \theta \left (\frac{\theta+1}{C\theta} \right )^{\frac{\theta}{\theta+1}}r^{-\frac{\tau - q\theta + \theta}{\theta+1}}\right ]\frac{C^\theta}{\theta} r^{\tau - q \theta -1} g(v(r)).
\end{equation}
Since $\tau > \max \{\theta (q-1), 0\}$, we can fix a large positive constant $R_0 > \tilde r$ such that
\begin{equation}\label{3.13}
q \theta \left (\frac{\theta+1}{C\theta} \right )^{\frac{\theta}{\theta+1}}r^{-\frac{\tau - q\theta + \theta}{\theta+1}} < \frac{1}{2}, \quad r>R_0.
\end{equation}
Then we have from \eqref{3.12} and \eqref{3.13} that 
\begin{equation}\label{3.14}
\frac{q}{r}\left( v^{\prime}(r)\right)^{\theta}< \frac{C^\theta}{2\theta} r^{\tau - q \theta -1} g(v(r)),\quad r>R_0.
\end{equation}
Inserting \eqref{3.14} into \eqref{3.1}, we then get \eqref{claim}.
\end{proof}

\end{document}
